\chapter*{{\color{oiB}Prakata}}
\markboth{\MakeUppercase{Prakata}}{}
%\chaptertext{}
%\sectiontext{}

\noindent%
OpenIntro Statistics mencakup materi mata kuliah statistika tingkat pengantar
dan menyajikan pengantar statistika terapan yang ketat secara ilmiah,
jelas, ringkas, dan mudah dipahami.
Buku ini ditulis terutama untuk jenjang sarjana,
tetapi juga banyak digunakan di sekolah menengah dan program pascasarjana.
\vspace{3mm}

Selain membangun landasan pemikiran dan metode statistika,
kami berharap pembaca memperoleh tiga gagasan berikut
dari buku ini.\vspace{-1mm}
\begin{itemize}
\setlength{\itemsep}{0mm}
\item
    Statistika adalah bidang terapan dengan cakupan
    penerapan praktis yang luas.
\item
    Anda tidak harus menjadi ahli matematika untuk belajar
    dari data nyata yang menarik.
\item
    Data sering kali tidak rapi, dan alat statistika tidak sempurna.
    Namun, jika memahami kekuatan dan keterbatasan alat-alat tersebut,
    Anda dapat menggunakannya untuk mempelajari dunia.
\end{itemize}


%\subsection*{Is this a data science book?}
%
%\noindent%
%Short answer: yes.
%Long answer: it depends what you mean by \term{data science},
%since two types of data scientists have emerged.
%\vspace{3mm}
%
%\noindent%
%Type~A data scientists focus on \emph{analysis},
%such as exploratory data analysis, inference,
%model building, and other related topics.
%Type~B data scientists focus on \emph{building},
%typically in the form of machine learning models
%or other systems.
%As you might expect, these two types share many skills,
%though their main focuses differ.
%This book focuses on skills most commonly used by
%Type~A data scientists.
%For more thoughts, please check out the following page:
%\begin{center}
%\oiRedirect{data_science_types}{{\color{red}BROKEN}}
%\end{center}
%\vspace{3mm}
%
%\noindent%


\subsection*{{\color{oiB}Ikhtisar buku ajar}}

\noindent%
Bab-bab dalam buku ini adalah sebagai berikut:%\vspace{2mm}
\begin{description}
\setlength{\itemsep}{0mm}
\item[1. Pengantar data.]
    Struktur data, variabel,
    dan teknik dasar pengumpulan data.
\item[2. Meringkas data.]
    Ringkasan data, grafik,
    dan pengantar singkat inferensi menggunakan pengacakan.
\item[3. Probabilitas.]
    Prinsip-prinsip dasar probabilitas.
    %This chapter is not required for the later chapters.
\item[4. Distribusi variabel acak.]
    Model normal dan distribusi penting lainnya.
\item[5. Dasar-dasar inferensi.]
    %Introduction to uncertainty in point estimates,
    %confidence intervals, and hypothesis tests.
    Gagasan umum statistika inferensial dalam konteks
    pendugaan proporsi populasi.
\item[6. Inferensi untuk data kategoris.]
    Inferensi untuk proporsi dan tabel dengan menggunakan distribusi
    normal dan khi-kuadrat.
\item[7. Inferensi untuk data numerik.]
    Inferensi untuk rata-rata satu atau dua sampel dengan menggunakan
    distribusi~\mbox{$t$},
    daya statistik untuk membandingkan dua kelompok,
    serta perbandingan banyak rata-rata
    menggunakan ANOVA.
\item[8. Pengantar regresi linear.]
    Regresi untuk hasil numerik dengan satu variabel prediktor.
    Sebagian besar bab ini dapat dipelajari setelah
    Bab~\ref{introductionToData}.
\item[9. Regresi berganda dan logistik.]
    Regresi untuk data numerik dan kategoris
    dengan menggunakan banyak prediktor. %for an accelerated course.
\end{description}


%\newpage

\noindent%
Struktur \emph{OpenIntro Statistics} memberi keleluasaan
dalam memilih dan mengurutkan topik.
Jika tujuan utama Anda adalah mempelajari regresi berganda
(Bab~\ref{ch_regr_mult_and_log})
secepat mungkin, prasyarat yang ideal adalah sebagai berikut:
\begin{itemize}
\setlength{\itemsep}{0mm}
\item Bab~\ref{ch_intro_to_data} dan
    Bagian~\ref{numericalData}--\ref{caseStudyMalariaVaccine}
    untuk memperoleh
    pengantar yang kokoh tentang struktur data dan ringkasan
    statistika yang digunakan di seluruh buku.
\item Bagian~\ref{normalDist}
    untuk memperoleh pemahaman yang kokoh tentang distribusi normal.
\item Bab~\ref{ch_foundations_for_inf}
    untuk membangun perangkat inti inferensi.
%\item Section~\ref{oneSampleMeansWithTDistribution}
%    and Chapter~\ref{ch_regr_simple_linear}
%    provide required for multiple regression with a numerical
%    outcome.
%    For the remaining chapters, they could be tackled in
%    almost any order, with the exception that
%    
%    
%    come before Chapter~\ref{ch_regr_mult_and_log}.
\item Bagian~\ref{oneSampleMeansWithTDistribution}
    untuk membangun landasan pemahaman tentang distribusi~$t$.
\item Bab~\ref{ch_regr_simple_linear}
    untuk membangun gagasan dan prinsip regresi
    dengan satu prediktor.
%    introduce the 
%    which introduces the $t$-distribution, should come before
%    Section~\ref{oneSampleMeansWithTDistribution}
%Chapters~\ref{ch_inference_for_props}-\ref{ch_regr_mult_and_log},
%    could be tackled in
%    almost any order, with the exception that
%    Section~\ref{oneSampleMeansWithTDistribution}
%    and Chapter~\ref{ch_regr_simple_linear}
%    come before Chapter~\ref{ch_regr_mult_and_log}.
%\item Sections~\ref{ch_inference_for_props}
%    and~\ref{} are recommended before logistic regression.
\end{itemize}
%One conspicuously missing topic from the list above is the
%chapter on Probability.
%While useful for a deeper understanding of the calculations,
%especially for anyone looking to take a second course in
%statistics, it is not required reading when the focus is on
%applied data analysis.


\subsection*{{\color{oiB}Contoh dan latihan}}
%, and appendices}

\noindent%
Contoh disediakan untuk membangun pemahaman tentang cara
menerapkan metode.

\begin{examplewrap}
\begin{nexample}{Ini adalah sebuah contoh.
    Jika pertanyaan diajukan di sini, di mana jawabannya dapat ditemukan?}
  Jawabannya dapat ditemukan di sini, pada bagian penyelesaian
  dalam contoh tersebut!
\end{nexample}
\end{examplewrap}

\noindent%
Jika kami menilai bahwa pembaca sudah siap mencoba menentukan
sendiri penyelesaian suatu contoh, kami menyajikannya sebagai Latihan Terbimbing.

\begin{exercisewrap}
\begin{nexercise}
Pembaca dapat memeriksa atau mempelajari jawaban setiap soal Latihan Terbimbing
dengan membaca penyelesaian lengkap dalam catatan kaki.\footnotemark{}
%Readers are strongly encouraged to attempt these practice problems.
\end{nexercise}
\end{exercisewrap}
\footnotetext{Soal Latihan Terbimbing dirancang untuk mengembangkan
  cara berpikir Anda. Anda dapat memeriksa jawaban dengan membaca
  penyelesaian dalam catatan kaki pada setiap Latihan Terbimbing.}

\noindent%
Latihan juga disediakan pada akhir setiap bagian,
bersama latihan ulasan pada akhir setiap bab.
Penyelesaian untuk latihan bernomor ganjil diberikan dalam
Lampiran~\ref{eoceSolutions}.
%Probability tables for the normal, $t$,
%and chi-square distributions are in
%Appendix~\ref{distributionTables}.


\subsection*{{\color{oiB}Sumber daya tambahan}}

Ikhtisar video, salindia, laboratorium perangkat lunak statistika,
himpunan data yang digunakan dalam buku ajar,
dan banyak sumber daya lainnya tersedia dengan mudah di\\[-5mm]
\begin{center}
\oiRedirect{os}
    {\color{black}\textbf{openintro.org/os}}
\end{center}
%Data sets for this textbook are available on the website
%and in a companion R package.\footnote{Diez DM,
%    Barr CD, \c{C}etinkaya-Rundel M. 2015.
%    \texttt{openintro}: OpenIntro data sets and supplement
%    functions.
%    \oiRedirect{textbook-github_openintro}
%        {github.com/OpenIntroOrg/openintro-r-package}.}
%All of these resources are free and may be used with
%or without this textbook as a companion.
Kami juga telah mempermudah akses ke data dalam buku ini
dengan menambahkan Lampiran~\ref{appendix_data},
yang memberikan informasi tambahan untuk setiap himpunan data
dalam teks utama dan merupakan bagian baru pada Edisi Keempat.
Panduan daring untuk setiap himpunan data tersebut juga tersedia di
\oiRedirect{data}
    {\color{black}\textbf{openintro.org/data}}
dan melalui
\oiRedirect{textbook-github_openintro}
    {paket pendamping R}.
% Official:
% http://www.openintro.org/package/openintro
% Currently redirect it to:
% http://openintrostat.github.io/openintro-r-package/
\vspace{3mm}

\noindent%
Kami menghargai semua umpan balik serta laporan kesalahan ketik
yang disampaikan melalui situs web.
Tautan singkat untuk melaporkan kesalahan ketik baru atau meninjau
kesalahan ketik yang telah diketahui adalah
\oiRedirect{os_typos}
    {\color{black}\textbf{openintro.org/os/typos}}. \vspace{3mm}

\noindent%
Bagi pembaca yang berfokus pada statistika di tingkat sekolah menengah,
pertimbangkan
\oiRedirect{textbook-books}
    {\emph{Advanced High School Statistics}},
yakni versi \emph{OpenIntro Statistics} yang telah
disesuaikan secara ekstensif oleh \oiRedirect{people}{Leah Dorazio}
untuk pelajaran sekolah menengah dan
AP\textsuperscript{\textregistered} Statistics.


\subsection*{{\color{oiB}Ucapan terima kasih}}
Proyek ini tidak akan terwujud tanpa semangat dan
dedikasi banyak orang selain mereka yang tercantum sebagai penulis.
Para penulis mengucapkan terima kasih kepada
\oiRedirect{textbook-openintro_about}{Staf OpenIntro}
atas keterlibatan dan kontribusi mereka yang berkelanjutan.
Kami~juga sangat berterima kasih kepada ratusan mahasiswa, siswa,
dan pengajar yang telah memberikan umpan balik berharga
sejak kami pertama kali mengunggah materi buku ini pada~2009. \vspace{3mm}

\noindent%
Kami juga berterima kasih kepada banyak pengajar yang membantu meninjau
edisi ini, termasuk
Laura Acion,
\oiRedirect{matthew_e_aiello-lammens}
    {Matthew E. Aiello-Lammens},
\oiRedirect{jonathan_akin}{Jonathan Akin},
Stacey C. Behrensmeyer,
Juan Gomez,
Jo Hardin,
\oiRedirect{nicholas_horton}{Nicholas Horton},
\oiRedirect{danish_khan}{Danish Khan},
\oiRedirect{peter_hm_klaren}{Peter H.M. Klaren},
Jesse Mostipak,
Jon C. New,
Mario Orsi,
Steve Phelps,
dan David Rockoff.
Kami menghargai semua umpan balik mereka, yang membantu
menyempurnakan teks ini secara berarti dan meningkatkan
kualitas buku ini secara substansial.
